\section{Conclusion}
\label{sec:conclusion}

We began with a paradox: zero-shot full-model SWA conversion of pre-trained transformers fails catastrophically, yet individual Llama-2-7B layers exhibit attention weight concentrated in fewer than 10\% of tokens---a universal pattern confirmed across 200 diverse evaluation examples without a single exception. If many layers are already operating with effectively local attention, why does converting them all at once cause collapse?

The answer, we argue, lies in selective identification. Not all layers are equally amenable to SWA conversion, and the difference between safe and unsafe conversion may be measurable---via per-layer attention entropy computed on a small calibration set. This paper establishes the empirical foundation for that argument.

\subsection*{Summary}

We addressed the problem of identifying which Llama-2-7B layers can be safely converted to sliding window attention without fine-tuning. Our approach uses head-mean attention entropy computed over 100 calibration sequences as a layer characterization signal. Our confirmed contributions are:

\textbf{(1)} Llama-2-7B exhibits near-universal within-layer attention concentration within the WikiText-103 evaluation domain (Gini $= 0.6829$, top-10\% token share $= 71.72\%$, both criteria satisfied by 100\% of 200 evaluation examples). This structural characterization---the first systematic application of per-layer head-mean entropy pooling to characterize attention concentration in Llama-2-7B for SWA layer selection---establishes that the model's attention is far from uniform global integration.

\textbf{(2)} Head-mean pooling is substantially more stable than head-max for capturing layer-level concentration (Gini 0.681 vs.\ 0.466, a 32\% relative difference). This is a methodological finding that applies beyond this work: any attention analysis method using per-layer entropy as a signal must justify its aggregation choice.

\textbf{(3)} Per-layer entropy rankings are stable across independent calibration subsets (Spearman $\rho \geq 0.8$ confirmed across all pairs), confirming that the entropy criterion produces deterministic layer selections in practice, independent of which 100 sequences are used for calibration.

Together, these findings validate the prerequisite for zero-shot selective SWA conversion: a stable, calibration-only layer characterization criterion exists. Whether this criterion successfully identifies SWA-safe layers---whether converting the 4 highest-entropy layers to SWA($w=512$) preserves WikiText-103 perplexity within 2 points---remains the central open question, pending h-e2 execution.

\textit{Note: Contribution (4)---the entropy-guided selective SWA framework---is a complete design and implementation artifact; its experimental validation (h-e2) is pending and not reported in this paper.}

\subsection*{Future Directions}

\textbf{From pending experiments (highest priority):} The immediate next step is executing h-e2 (entropy-guided $k=4$ SWA conversion on WikiText-103 test set) and h-m1 (entropy vs.\ random vs.\ last-$k$ comparison). Regardless of outcome, h-e2 provides publishable findings: confirmation validates the full framework; refutation characterizes the zero-shot SWA feasibility boundary in Llama-2-7B. The exploratory QA F1 directional signal (0.43 pp vs.\ 0.67 pp random, non-significant at $p = 0.4507$) motivates the proper h-m1 experiment.

\textbf{From unverified assumptions:} The residual stream compensation mechanism---that 28 remaining full-attention layers compensate for 4 converted SWA layers---is theoretically motivated but empirically unconfirmed. Depth-position analysis in h-m2 (comparing $k=4$ vs.\ $k=8$ degradation) can test whether early-layer conversion causes more degradation than late-layer conversion.

\textbf{From scope extensions:} Cross-domain calibration stability is the most impactful near-term extension: re-running entropy scoring with SST-2 or code-domain calibration sequences would test whether the layer rankings are consistent structural properties or domain-specific phenomena. Architecture generalization to Llama-3-8B (grouped-query attention, requiring adapted pooling) is a natural next step after Llama-2-7B validation.

Our findings suggest that pre-trained LLMs contain latent structural heterogeneity---some layers are already operating locally---and that this heterogeneity can be measured cheaply and reliably within a single evaluation domain. We hope this work encourages the community to look beyond holistic model modification and toward layer-level characterization as a principled basis for efficient, training-free inference optimization.
